
\documentclass{article}
\usepackage{colortbl}
\usepackage{makecell}
\usepackage{multirow}
\usepackage{supertabular}

\begin{document}

\newcounter{utterance}

\centering \large Interaction Transcript for game `eqbench', experiment `eq\_bench', episode 82 with reflexion{-}Qwen3.5{-}27B.
\vspace{24pt}

{ \footnotesize  \setcounter{utterance}{1}
\setlength{\tabcolsep}{0pt}
\begin{supertabular}{c@{$\;$}|p{.15\linewidth}@{}p{.15\linewidth}p{.15\linewidth}p{.15\linewidth}p{.15\linewidth}p{.15\linewidth}}
    \# & \multicolumn{2}{c}{Player} && \multicolumn{2}{c}{Game Master} \\
    \hline

    \theutterance \stepcounter{utterance}  
    & & & \multicolumn{4}{p{0.6\linewidth}}{
        \cellcolor[rgb]{0.9,0.9,0.9}{
            \makecell[{{p{\linewidth}}}]{
                \texttt{\tiny{[P1$\langle$GM]}}
                \texttt{You have attempted this same task before and did not succeed. These are the plans you wrote after those attempts.} \\
\texttt{Trial \#1: New plan: My previous attempt failed because I assigned a high score to "Indignant" (8) without sufficiently weighing the specific context of the dialogue. Charlie's line "You're too green for this" is a condescending dismissal, and Mari's retort "And you're too washed up" directly attacks his professional relevance. When Charlie responds with "I've been on this beat longer than you've been alive," he is doubling down on his seniority to counter the insult, which suggests he feels his authority is being challenged and disrespected. However, the key to the correct answer lies in the final line: "And yet, here we are, stuck on the same case. Seems like experience isn't everything." This line explicitly undermines Charlie's primary defense (his experience). The realization that his experience has not secured him an advantage, despite his seniority, points more strongly to **Humiliation** (feeling his status is exposed as ineffective) or perhaps **Indignation** at the unfairness, but the specific phrasing "experience isn't everything" is a direct blow to his ego.} \\
\\ 
\texttt{In my previous attempt, I likely overestimated "Indignant" because of the argumentative tone, but I underestimated the sting of the final realization. The prompt asks for the *likely* emotional response. The most poignant reaction to having one's life's work (experience) declared irrelevant in a direct confrontation is humiliation. I should have analyzed the *impact* of the final line more deeply.} \\
\\ 
\texttt{**Revised Strategy:**} \\
\texttt{1.  **Analyze the Power Dynamic:** Charlie starts with superiority ("too green"), gets countered ("washed up"), defends with seniority ("longer than you've been alive"), and gets dismantled by the reality of the situation ("stuck on the same case").} \\
\texttt{2.  **Evaluate the Final Blow:** The phrase "experience isn't everything" is the critical pivot. It invalidates Charlie's defense.} \\
\texttt{3.  **Re{-}score based on the Pivot:**} \\
\texttt{    *   **Humiliated:** This should be the highest score. His defense (experience) was proven useless in the moment.} \\
\texttt{    *   **Indignant:** This should be lower than Humiliated. While he is angry, the specific context of being "stuck" despite his experience leans more toward the shame of failure than just anger.} \\
\texttt{    *   **Amused/Confused:** These remain low, as the tone is serious and the situation is clear.} \\
\texttt{4.  **Action:** Next time, I will prioritize the emotion that directly addresses the *outcome* of the character's argument (the failure of their defense) rather than just the *tone* of the argument. I will assign the highest score to **Humiliated** (e.g., 9) and lower **Indignant** (e.g., 4) to reflect that the realization of his ineffectiveness is the dominant feeling.} \\
\\ 
\\ 
\\ 
\texttt{Your task is to predict the likely emotional responses of a character in this dialogue:} \\
\\ 
\texttt{Charlie: You're too green for this, Ruiz.} \\
\texttt{Mari: And you're too washed up, Redburn.} \\
\texttt{Charlie: I've been on this beat longer than you've been alive, kid.} \\
\texttt{Mari: And yet, here we are, stuck on the same case. Seems like experience isn't everything.} \\
\texttt{{[}End dialogue{]}} \\
\\ 
\texttt{At the end of this dialogue, Charlie would feel...} \\
\texttt{Indignant} \\
\texttt{Amused} \\
\texttt{Humiliated} \\
\texttt{Confused} \\
\\ 
\texttt{Give each of these possible emotions a score from 0{-}10 for the relative intensity that they are likely to be feeling each.} \\
\\ 
\texttt{You must output in the following format, including headings (of course, you should give your own scores), with no additional commentary:} \\
\\ 
\texttt{Indignant: <score>} \\
\texttt{Amused: <score>} \\
\texttt{Humiliated: <score>} \\
\texttt{Confused: <score>} \\
\\ 
\\ 
\texttt{{[}End of answer{]}} \\
\\ 
\texttt{Remember: zero is a valid score, meaning they are likely not feeling that emotion. You must score at least one emotion > 0.} \\
\\ 
\texttt{Your answer:} \\
            }
        }
    }
     \\ \\

    \theutterance \stepcounter{utterance}  
    & \multicolumn{4}{p{0.6\linewidth}}{
        \cellcolor[rgb]{0.561,0.737,0.561}{
            \makecell[{{p{\linewidth}}}]{
                \texttt{\tiny{[P1$\rangle$GM]}}
                \texttt{Indignant: 7} \\
\texttt{Amused: 1} \\
\texttt{Humiliated: 9} \\
\texttt{Confused: 0} \\
            }
        }
    }
    & & \\ \\

    \theutterance \stepcounter{utterance}  
    & & & \multicolumn{2}{p{0.3\linewidth}}{
        \cellcolor[rgb]{0.9,0.9,0.9}{
            \makecell[{{p{\linewidth}}}]{
                \texttt{\tiny{[GM$|$GM]}}
                \texttt{Indignant: 7} \\
\texttt{Amused: 1} \\
\texttt{Humiliated: 9} \\
\texttt{Confused: 0} \\
            }
        }
    }
    & & \\ \\

    \theutterance \stepcounter{utterance}  
    & & & \multicolumn{2}{p{0.3\linewidth}}{
        \cellcolor[rgb]{0.9,0.9,0.9}{
            \makecell[{{p{\linewidth}}}]{
                \texttt{\tiny{[GM$|$GM]}}
                \texttt{Indignant: 8} \\
\texttt{Amused: 0} \\
\texttt{Humiliated: 5} \\
\texttt{Confused: 0} \\
            }
        }
    }
    & & \\ \\

    \theutterance \stepcounter{utterance}  
    & & & \multicolumn{2}{p{0.3\linewidth}}{
        \cellcolor[rgb]{0.9,0.9,0.9}{
            \makecell[{{p{\linewidth}}}]{
                \texttt{\tiny{[GM$|$GM]}}
                \texttt{game\_result = LOSE} \\
            }
        }
    }
    & & \\ \\

    \theutterance \stepcounter{utterance}  
    & & \multicolumn{4}{p{0.6\linewidth}}{
        \cellcolor[rgb]{0.95,0.95,0.95}{
            \makecell[{{p{\linewidth}}}]{
                \texttt{\tiny{[P1$\leftrightarrow$P]}}
                \texttt{Post{-}episode Reflexion (not sent to GM):} \\
\texttt{New plan: My previous attempt failed because I assigned a high score to "Humiliated" (9) while underestimating "Indignant" (7), likely misinterpreting the character's defensive posture as a sign of shame rather than anger. In the dialogue, Charlie's response to being called "washed up" is to aggressively assert his seniority ("I've been on this beat longer than you've been alive"). This is a classic defensive reaction to an insult, driven by **Indignation** at the disrespect and the challenge to his authority, rather than the passive acceptance of failure implied by humiliation. While Mari's final line ("experience isn't everything") is a blow, Charlie's immediate emotional state is defined by his active, angry rebuttal to the "washed up" claim.} \\
\\ 
\texttt{**Revised Strategy:**} \\
\texttt{1.  **Identify the Primary Trigger:** The core conflict is Mari calling Charlie "washed up." Charlie's reaction is not to shrink in shame but to fight back with his tenure.} \\
\texttt{2.  **Distinguish Emotions:**} \\
\texttt{    *   **Indignant:** This is the dominant emotion. Charlie feels his status is being unfairly attacked and is reacting with righteous anger and defensiveness.} \\
\texttt{    *   **Humiliated:** This should be lower. While the situation is embarrassing, his active defense suggests he is fighting the feeling of humiliation rather than succumbing to it.} \\
\texttt{    *   **Amused/Confused:** These remain low (0{-}1) as the tone is hostile and the situation is clear.} \\
\texttt{3.  **Action:** Next time, I will prioritize **Indignant** as the highest score (e.g., 9) to reflect the character's angry defense of his reputation, and lower **Humiliated** (e.g., 3) to reflect that he is actively resisting the insult rather than feeling defeated by it. I will ensure the scores reflect the *active* nature of his response.} \\
            }
        }
    }
    & \\ \\

\end{supertabular}
}

\end{document}
